\documentclass[10pt,a4paper]{amsart}
\usepackage[margin=19mm]{geometry}
\usepackage{amsmath,amssymb,mathtools}
\usepackage[scr=rsfs,cal=euler]{mathalfa}
\usepackage{xfrac}
\usepackage[unicode,hidelinks,bookmarks=false]{hyperref}
\newcommand{\ie}{i.e.\ }
\newcommand{\cf}{cf.\ }
\newcommand{\resp}{respectively\ }
\newcommand{\Subsection}{\subsection}
\providecommand{\nobreakdash}{\nobreak\hbox{-}}
\setlength{\emergencystretch}{2em}
\multlinegap0pt
% ---------------------------------------------------------------------------
% Editorial AMS wrapper prelude. The theorem environments and the mathematical macros below are
% transcribed from the paper's own preamble (cached-originals/source0.tex lines 1--44); the AMS amsart
% class, the named format packages of the pinned TeX Live runtime and this editorial formatting are not
% part of the author's text. No mathematics is changed.
% ---------------------------------------------------------------------------
\usepackage{cleveref}
\usepackage{cleveref}
\allowdisplaybreaks[4]
\newtheorem{corollary}{Corollary}
\newtheorem{definition}{Definition}
\newtheorem{lemma}{Lemma}
\newtheorem{problem}{Problem}
\newtheorem{remark}{Remark}
\newtheorem{theorem}{Theorem}
\newtheorem{prop}{Proposition}
\newcommand\adots{\mathinner{\mkern2mu%
	\raisebox{0.1em}{.}\mkern2mu\raisebox{0.4em}{.}%
	\mkern2mu\raisebox{0.7em}{.}\mkern1mu}}
\newcommand\Ao{A_\mathrm o}
\newcommand\calC{\mathcal C}
\newcommand\calR{\mathcal R}
\newcommand\faii{\varphi_\mathrm i}
\newcommand\faiibar[1]{\bar\varphi_\mathrm i^{[#1]}}
\newcommand\GAmma{\mathit\Gamma}
\newcommand\nsharp{n_\mathrm i^*}
\newcommand\nocstar{n_\mathrm{oc}^*}
\newcommand\Pio{\mathit\Omega_\mathrm o}
\newcommand\Qi{Q_\mathrm i}
\newcommand\Qoc{Q_\mathrm{oc}}
\newcommand\rank{\mathrm{rank}}
\newcommand\rref{\mathrm{rref}}
\newcommand\im{\mathrm{im}}
\newcommand\imr[1]{\mathrm{im}(#1^\mathrm T)}
\newcommand\R{\mathbf R}
\newcommand\T{\mathrm T}
\newcommand\PHi{\mathit\Phi}
\newcommand\PSi{\mathit\Psi}

\begin{document}
\title{Determination of the output injection of invertible linear systems for guaranteeing a regular relative degree}
\author{Shida Cao, Bin Zhou, Guang-Ren Duan}
\date{Source: arXiv:2609.20018v1; distributed under CC BY 4.0}
\maketitle
\noindent\emph{Source, version and licence.} \emph{Determination of the output injection of invertible linear systems for guaranteeing a regular relative degree}, by Shida Cao, Bin Zhou, Guang-Ren Duan. The source of this editable unit is the e-print of \textbf{version v1} of arXiv:2609.20018, https://arxiv.org/abs/2609.20018v1, whose material is distributed under the Creative Commons Attribution 4.0 International licence, http://creativecommons.org/licenses/by/4.0/, as recorded in the arXiv licence metadata for this paper and shown on that abstract page. The whole of the body below is version v1, the newest version of the paper. Full rights notice: licenses/arxiv-2609.20018-CC-BY-4.0.txt.\par\smallskip
\noindent\textit{Editorial scope.} Counted unit: original Theorem 1 of the article, the statement carried by the source environment \texttt{theorem} with source label \texttt{thm: (A,B,TyC) and (A+LC,B,TyC), Ty the same}, cached-originals/source0.tex lines 849--858, printed as \textbf{Theorem 1} exactly as in the article: for an invertible system with n-sharp equal to two and any output transformation matrix, the system (A,B,T\_yC) has a regular relative degree if and only if the output-injected system (A+LC,B,T\_yC) has a regular relative degree for every Luenberger gain L. It is delivered together with the authors' own complete proof as the single primary proof body, source0.tex lines 860--890 (31 lines): the \texttt{proof} environment that belongs to the statement and establishes it in full. No proof marker is added and no mathematics is written, repaired or re-derived; the authors' notation and the printed slips of the source are kept literal. Necessary complete context is retained uncounted, in source order: lines 273--275, the source \texttt{equation} with source label \texttt{eqn: xdot=Ax+Bu, y=Cx} of the article that the retained passages refer to; lines 281--301, the source \texttt{definition} with source label \texttt{def:rd} of the article that the retained passages refer to; lines 751--754, the source \texttt{lemma} with source label \texttt{thm: if nsharp=2 then n*=2} of the article that the retained passages refer to. The article's own numbering is reproduced by explicit counter settings: the theorem-like counter of the article is set before each retained passage, so the counted statement prints with the number the article gives it. No \texttt{\textbackslash cite} command occurs in any retained passage: the closure of the retained passages over references and citations was computed, it contains no citation at all, so this unit delivers no bibliography entry and the article's own bibliography data is neither spliced into the bundled source nor redistributed. The source's own class is \texttt{elsarticle}; the e-print ships no class or style file, so no publisher class is shipped, used or redistributed, and the wrapper is the AMS \texttt{amsart} class with the article's own macros and theorem declarations transcribed from its preamble. The retained passages contain no figure environment and no graphics-inclusion command, so no artwork is redistributed. Every declared body of this unit is an exact, unmodified span of source0.tex: no body is a bounded passage comparison, \texttt{omit} was not used anywhere, and consequently no fidelity receipt is asserted in delivery-evidence.json. The complete concatenated original source is bundled as sources/210827/source0.tex. Omitted from this editable unit and therefore absent from it are source0.tex lines 1--272; 276--280; 302--750; 755--848; 859--859; 891--2916. The AMS \texttt{amsart} wrapper, the named format packages of the pinned TeX Live runtime and this editorial source, licence and scope paragraph are editorial formatting only.\par\smallskip
\setcounter{section}{2}
\setcounter{equation}{0}
\setcounter{corollary}{0}
\setcounter{definition}{0}
\setcounter{lemma}{0}
\setcounter{problem}{0}
\setcounter{prop}{0}
\setcounter{remark}{0}
\setcounter{theorem}{0}
\begin{equation}\label{eqn: xdot=Ax+Bu, y=Cx}
\dot x=Ax+Bu,\ y=Cx.
\end{equation}

\setcounter{section}{2}
\setcounter{equation}{1}
\setcounter{corollary}{0}
\setcounter{definition}{0}
\setcounter{lemma}{0}
\setcounter{problem}{0}
\setcounter{prop}{0}
\setcounter{remark}{0}
\setcounter{theorem}{0}
\begin{definition}[regular relative degree]\label{def:rd}
Denote $C=\begin{bmatrix}c_1^\T&c_2^\T&\cdots&c_m^\T\end{bmatrix}^\T$
in system \eqref{eqn: xdot=Ax+Bu, y=Cx},
where $c_i\in\R^{1\times n}$ for $i=1,2,\dots,m$.
The ordered set $\{r_1,r_2,\dots,r_m\}$, where $1\le r_i\le n$,
is said to be the relative degree of system
\eqref{eqn: xdot=Ax+Bu, y=Cx} if for $i=1,2,\dots,m$, there holds
$c_iA^jB=0$ for $j=0,1,\dots,r_i-2$, $c_iA^{r_i-1}B\ne0$.
We say the relative degree $\{r_1,r_2,\dots,r_m\}$ is regular if
$\det(\GAmma(A,B,C))\ne0$, where
\[
\GAmma(A,B,C)=
\begin{bmatrix}
(c_1A^{r_1-1}B)^\T&(c_2A^{r_2-1}B)^\T&\cdots&(c_mA^{r_m-1}B)^\T
\end{bmatrix}^\T.
\]
Here, $\GAmma(A,B,C)\in\R^{m\times m}$ is referred to
as the coupling matrix.
Conversely, the relative degree $\{r_1,r_2,\dots,r_m\}$
is said to be irregular if $\det(\GAmma(A,B,C))=0$.
\end{definition}

\setcounter{section}{3}
\setcounter{equation}{0}
\setcounter{corollary}{0}
\setcounter{definition}{4}
\setcounter{lemma}{6}
\setcounter{problem}{1}
\setcounter{prop}{2}
\setcounter{remark}{4}
\setcounter{theorem}{0}
\begin{lemma}\label{thm: if nsharp=2 then n*=2}
For the system \eqref{eqn: xdot=Ax+Bu, y=Cx},
if $\nsharp=2$, then $\nocstar=2$.
\end{lemma}

\setcounter{section}{3}
\setcounter{equation}{3}
\setcounter{corollary}{0}
\setcounter{definition}{4}
\setcounter{lemma}{7}
\setcounter{problem}{1}
\setcounter{prop}{3}
\setcounter{remark}{4}
\setcounter{theorem}{0}
\begin{theorem}
\label{thm: (A,B,TyC) and (A+LC,B,TyC), Ty the same}
Consider the invertible system \eqref{eqn: xdot=Ax+Bu, y=Cx}
with $\nsharp=2$.
Let $T_y\in\R^{m\times m}_m$ be an output transformation matrix.
Then for any Luenberger gain $L\in\R^{n\times m}$,
the system $(A,B,T_yC)$ has a regular relative degree
if and only if
the system $(A+LC,B,T_yC)$ has a regular relative degree.
\end{theorem}

\setcounter{section}{3}
\setcounter{equation}{3}
\setcounter{corollary}{0}
\setcounter{definition}{4}
\setcounter{lemma}{7}
\setcounter{problem}{1}
\setcounter{prop}{3}
\setcounter{remark}{4}
\setcounter{theorem}{1}
\begin{proof}
We focus on proving the necessity, as the sufficiency follows by a similar argument and is thus omitted for brevity.

According to \Cref{thm: if nsharp=2 then n*=2}, the condition $\nsharp=2$ implies $\nocstar=2$, which yields
\[
\rank(T_yCB)<m, \ \rank\begin{bmatrix}T_yCB & T_yCAB\end{bmatrix}=m.
\]
Consequently, the relative degree of the system $(A,B,T_yC)$ must be either $\{1,\dots,1,2,\dots,2\}$ or $\{2,\dots,2\}$.
Suppose the relative degree is
\[
\{\underbrace{1,\dots,1}_{n_1\,\mathrm{times}}, \underbrace{2,\dots,2}_{n_2\,\mathrm{times}}\},
\text{ where $n_1$ may be $0$.}
\]
Without loss of generality, let $T_y$ be partitioned as $T_y=\begin{bmatrix}T_1\\ T_2\end{bmatrix}$, where $T_1\in\R^{n_1\times m}$ and $T_2\in\R^{n_2\times m}$.
According to the definition of the regular relative degree (\Cref{def:rd}), it holds that $T_2CB=0$, and the coupling matrix
$\GAmma(A,B,T_yC)=\begin{bmatrix}T_1CB\\ T_2CAB\end{bmatrix}$
is nonsingular. By noting that $T_yCB = \begin{bmatrix}T_1CB \\ 0\end{bmatrix}$ and partitioning $T_2CLT_y^{-1}$ as $\begin{bmatrix}T_3 & *\end{bmatrix}$ where $T_3\in\R^{n_2\times n_1}$, we obtain
\begin{align*}
T_2C(A+LC)B
&= T_2CAB + T_2CLCB
= T_2CAB + T_2CLT_y^{-1}T_yCB \\
&= T_2CAB + \begin{bmatrix}T_3 & *\end{bmatrix} \begin{bmatrix}T_1CB \\ 0\end{bmatrix}
= T_2CAB + T_3T_1CB.
\end{align*}
Since $\begin{bmatrix}T_1CB \\ T_2CAB\end{bmatrix}$ is nonsingular, it follows that each row of $T_2CAB + T_3T_1CB$ is non-zero.
Thus, the coupling matrix is
\[
\GAmma(A+LC,B,T_yC) = \begin{bmatrix}T_1CB \\ T_2C(A+LC)B\end{bmatrix} = \begin{bmatrix}I_{n_1} & \\ T_3 & I_{n_2}\end{bmatrix} \begin{bmatrix}T_1CB \\ T_2CAB\end{bmatrix},
\]
and it is also nonsingular. This implies that the system $(A+LC,B,T_yC)$ also possesses a regular relative degree, which completes the proof.
\end{proof}

\end{document}
